% Responsable: Integrante 2 (tarea 2.2 del plan)
\section{Especificaciones de diseño}\label{sec:especificaciones}
Se considera que el bloque frío funciona como una pequeña cámara para conservar muestras o vacunas. La Organización Mundial de la Salud establece para su almacenamiento el rango de $+2$ a $+8\,^\circ\mathrm C$ \cite{who_pqs}; la referencia de $5\,^\circ\mathrm C$ queda en el centro de ese rango. Para la medición se supone un sensor digital DS18B20, con resolución de $0.0625\,^\circ\mathrm C$ en 12 bits y exactitud de $\pm0.5\,^\circ\mathrm C$ entre $-10$ y $+85\,^\circ\mathrm C$ \cite{ds18b20}. Con este contexto se definen las especificaciones de la Tabla~\ref{tab:especificaciones}. Las especificaciones E1 a E7 se evalúan con el modelo lineal alrededor del punto de operación ($T_{b0}=5\,^\circ\mathrm C$, $I_0=1.19\,\mathrm A$); la E8 se evalúa con el modelo no lineal.

\begin{table*}[!tbp]\centering
\caption{Especificaciones de diseño del lazo de temperatura.}\label{tab:especificaciones}
{\tableformat
\begin{tabularx}{\textwidth}{c>{\raggedright\arraybackslash}p{3.4cm}>{\raggedright\arraybackslash}p{5.2cm}L}\toprule
Id. & Métrica & Especificación & Justificación\\\midrule
E1 & Error en estado estable & $|e_{ss}|\le0.1\,\mathrm K$ para escalones de hasta $3\,\mathrm K$ & Equivale a menos de dos pasos de resolución del sensor \cite{ds18b20}. Un cambio de $3\,\mathrm K$ cubre toda la mitad del rango de $2$ a $8\,^\circ\mathrm C$ \cite{who_pqs}.\\
E2 & Sobrepaso & $M_p\le10\,\%$ & Para un escalón de $1\,\mathrm K$, el sobrepaso queda del mismo orden que la tolerancia de E1.\\
E3 & Tiempo de asentamiento (2\,\%) & $t_s\le180\,\mathrm s$ & Aproximadamente la mitad del tiempo de lazo abierto ($373\,\mathrm s$), para recuperar la temperatura después de abrir la cámara o introducir una muestra.\\
E4 & Esfuerzo de control en pequeña señal & $\Delta I_{\max}\le0.81\,\mathrm A$ ante un escalón de $-1\,\mathrm K$ & Con $I_0=1.19\,\mathrm A$, la corriente total no alcanza el límite de $2\,\mathrm A$ y el lazo trabaja sin saturar. La Tarea~1 mostró un error del modelo lineal cercano al $5\,\%$ con $|\Delta I|=1\,\mathrm A$.\\
E5 & Rango de la corriente & $0\le I\le2.0\,\mathrm A$ & Por debajo del mínimo de la curva estática ($2.42\,\mathrm A$, Figura~\ref{fig:limite}); la fuente no invierte la corriente.\\
E6 & Rechazo de perturbación & Escalón de $+5\,\mathrm K$ en $T_{\rm amb}$ o carga de $2\,\mathrm W$ en el bloque: pico $\le0.5\,\mathrm K$, regreso a $\pm0.1\,\mathrm K$ en $\le400\,\mathrm s$ y error final nulo & Un cambio de $25$ a $30\,^\circ\mathrm C$ en el ambiente es esperable durante el día, y la carga de $2\,\mathrm W$ representa introducir una muestra. En lazo abierto producirían $+4.6$ y $+5.0\,\mathrm K$, lo que saca al bloque del rango de $2$ a $8\,^\circ\mathrm C$.\\
E7 & Estabilidad relativa & Margen de fase $\ge45^\circ$ & Las masas del bloque y del disipador se marcaron como provisionales en la Tarea~1; el margen absorbe errores en esos parámetros.\\
E8 & Arranque desde $25\,^\circ\mathrm C$ (no lineal) & $T_b\ge2\,^\circ\mathrm C$ en todo momento y $|T_b-5\,^\circ\mathrm C|\le0.5\,\mathrm K$ antes de $300\,\mathrm s$ & No se aceptan temperaturas por debajo de $+2\,^\circ\mathrm C$ en el almacenamiento \cite{who_pqs}. Se exige además que el arranque sea más rápido que el de lazo abierto.\\\bottomrule
\end{tabularx}}
\par\smallskip\footnotesize Fuente: elaboración propia.
\end{table*}

Con las ecuaciones~\eqref{eq:essP} y \eqref{eq:essD} se ve que el control P no puede cumplir todas las especificaciones al mismo tiempo. La E1 exige $K_p\ge2.81\,\mathrm{A/K}$ y la E6 exige $K_p\ge4.35\,\mathrm{A/K}$ por la temperatura ambiente y $K_p\ge4.76\,\mathrm{A/K}$ por la carga de $2\,\mathrm W$, pero la E4 limita la ganancia a $K_p\le0.81\,\mathrm{A/K}$. Ningún valor de $K_p$ satisface las tres condiciones, por lo que se necesita un compensador.
